% Postnikov truncations as joins of sheaf and homotopy equivalences.
% All finite claims verified by scripts/comparison/sheaf_sameness*.py.
\documentclass[11pt]{amsart}
\usepackage{amsmath,amssymb}
\usepackage{stmaryrd}
\setlength{\emergencystretch}{3em}
\tolerance=600
\usepackage{tikz-cd}
\usepackage{booktabs}
\usepackage[hidelinks]{hyperref}

\newtheorem{thm}{Theorem}[section]
\newtheorem{lem}[thm]{Lemma}
\newtheorem{cor}[thm]{Corollary}
\newtheorem{prop}[thm]{Proposition}
\theoremstyle{definition}
\newtheorem{defi}[thm]{Definition}
\newtheorem{rem}[thm]{Remark}
\newtheorem{ex}[thm]{Example}

\newcommand{\Same}{\mathrm{Same}}
\newcommand{\Sub}{\mathrm{Sub}}
\newcommand{\Top}{\mathrm{Top}}
\newcommand{\Mor}{\mathrm{Mor}}
\newcommand{\Iso}{\mathrm{Iso}}
\newcommand{\Presh}{\mathrm{Presh}}
\newcommand{\Sh}{\mathrm{Sh}}
\newcommand{\Down}{\mathrm{Down}}
\newcommand{\cosk}{\mathrm{cosk}}
\newcommand{\sk}{\mathrm{sk}}
\newcommand{\Wtest}{\mathsf{W}_{\mathrm{test}}}
\newcommand{\Wtype}{\mathsf{W}_{\mathrm{type}}}
\newcommand{\sset}{\mathbf{sSet}}
\newcommand{\Kan}{\mathbf{Kan}}
\newcommand{\cube}{\square}
\newcommand{\Idem}{\mathrm{Idl}^{\,2}_{\mathrm{idem}}}

\begin{document}

\title[Postnikov truncations as joins]{Postnikov truncations as joins of
  sheaf and homotopy equivalences}
\author{Ryota Shibusawa}
\thanks{Corresponding author. E-mail: \texttt{r-shibusawa@daiichi-koudai.ac.jp}}
\address{Daiichi Institute of Technology}
\email{r-shibusawa@daiichi-koudai.ac.jp}
\keywords{Grothendieck topology, subtopos, sheafification, localization, weak
  equivalence, two-out-of-three, idempotent ideal, semilattice, coskeleton,
  Postnikov truncation, Kan complex, modal logic}
\subjclass[2020]{18F10 (primary), 18F20, 18A99, 06A12, 20M12, 55U10 (secondary)}

\begin{abstract}
In the complete lattice $\Same(\sset)$ of all two-out-of-three classes of
morphisms of simplicial sets, let $\Wtest$ be the weak homotopy equivalences
and $W_{D_k}$ the maps inverted by the $k$-coskeleton, i.e.\ the sheaf
equivalences of the $k$-th coskeletal subtopos. We prove that for every
$n\ge-1$ the join $\Wtest\vee W_{D_{n+1}}$ is exactly the class of
$n$-truncated weak equivalences: the Postnikov truncation is the lattice join
of sheaf sameness and homotopy sameness, with no fibrancy hypothesis, and
$\Wtest=\bigcap_n(\Wtest\vee W_{D_n})$ (a sheaf-theoretic Whitehead theorem).
The proof is an instance of a general join theorem: if an endofunctor $P$ is
connected to the identity by a natural transformation with components in
$U\vee V$ and sends $U$ and $V$ into a sameness $W\subseteq U\vee V$, then
$U\vee V=P^{*}(W)$; for two reflectors $L_1,L_2$ this gives
$W_{L_1}\vee W_{L_2}=W_{L_2L_1}$ iff $L_2L_1$ inverts $W_{L_2}$. To set this
up we show that for any Grothendieck topology $J$ on a small category $S$ the
class $W_J$ of maps inverted by sheafification is an element of
$\Same(\Presh(S))$, that $J\mapsto W_J$ is an order-embedding whose image
consists of the saturated samenesses inverted by a left-exact reflector, and
that, within the embedded sheaf samenesses, joins in the lattice of
topologies are obtained by saturating joins in the sameness lattice:
$W_{J\vee J'}=\mathrm{Sat}(W_J\vee W_{J'})$. We also give a short direct proof, in this language, of the
finite-sieve form of the Kelly--Lawvere correspondence --- the Grothendieck
topologies on a category with finitely many sieves on each object are its
idempotent two-sided ideals --- and re-derive from it the classical facts that
the subtoposes of simplicial and cubical sets are the coskeletal ones; the
dictionary is used to locate the sheaf chain $W_{D_k}$ inside
$\Same(\sset)$. All finite claims are verified by exhaustive computation from
the definitions.
\end{abstract}

\maketitle

\section{Introduction}

Two theories may agree on what an object \emph{is} and disagree on when two
objects are \emph{the same}. Homotopy theory identifies simplicial sets along
weak equivalences; sheaf theory identifies presheaves that have the same
sheafification. Both notions of sameness are classes of morphisms satisfying
two-out-of-three, and the companion papers \cite{SameA,SameB} proposed to
study all such classes at once: for a category $C$, a \emph{sameness} is a
class $W$ of morphisms containing the identities, closed under composition and
satisfying two-out-of-three, and $\Same(C)$ is the complete lattice of
samenesses under inclusion. The present paper compares the sheaf-theoretic
and the homotopical samenesses of simplicial sets \emph{inside} this lattice.

\medskip\noindent\textbf{Main theorem.} Let $\Wtest\in\Same(\sset)$ be the
weak homotopy equivalences and, for $k\ge0$, let $W_{D_k}$ be the class of
maps $f$ with $\cosk_kf$ an isomorphism --- equivalently, bijective in
dimensions $\le k$ --- which is the class of maps inverted by sheafification
for the $k$-th coskeletal subtopos. For $n\ge-1$ let $T_n$ be the
$n$-truncated weak equivalences (maps inducing isomorphisms on $\pi_i$ for
$i\le n$ at all base points; for $n=-1$, maps whose source and target are
both empty or both non-empty).
Then (Theorem~\ref{thm:join} and Corollary~\ref{cor:whitehead})
\[
\Wtest\vee W_{D_{n+1}}=T_n\quad(n\ge-1),
\]
\[
\Wtest=\bigcap_{n\ge0}\bigl(\Wtest\vee W_{D_n}\bigr),\qquad
\Iso=\bigcap_{n\ge0}W_{D_n}.
\]
That the coskeleta of a \emph{Kan complex} model its Postnikov sections is
classical; the point here is the statement about \emph{classes of maps}: the
Postnikov truncation, which at the $\infty$-level is the $n$-truncation
modality, is the lattice join of the sheaf sameness of the coskeletal
topology with the homotopy sameness, on all simplicial sets, and the homotopy
sameness is the meet of these joins. In words,
\[
\text{sheaf localization}\ \vee\ \text{homotopy localization}\ =\ \text{Postnikov localization}.
\]

\medskip\noindent\textbf{The general join theorem.} The proof is an instance
of a statement about arbitrary samenesses (Theorem~\ref{thm:generaljoin}):
if an endofunctor $P$ of $C$ receives a natural transformation from the
identity whose components lie in $U\vee V$, and $P$ sends $U$ and $V$ into a
sameness $W\subseteq U\vee V$, then $U\vee V=P^*(W)$, the class of maps sent
by $P$ into $W$. For two reflectors $L_1,L_2$ onto full subcategories, with
$W_{L_i}$ the maps they invert, this specialises to
$W_{L_1}\vee W_{L_2}=W_{L_2L_1}$ iff $L_2L_1$ inverts $W_{L_2}$
(Corollary~\ref{cor:tworeflectors}); for two Grothendieck topologies $J,J'$
it shows that $W_{J\vee J'}$ is the saturation of $W_J\vee W_{J'}$
(Proposition~\ref{prop:topjoin}). The main theorem is the case $U=\Wtest$,
$V=W_{D_{n+1}}$, $P=\cosk_{n+1}\circ\mathrm{Ex}^\infty$, $W=\Wtest$.

\medskip\noindent\textbf{Sheaf sameness.} To place sheaf theory in the
lattice we show (Theorem~\ref{thm:embed}) that for every Grothendieck
topology $J$ on a small category $S$ the class $W_J$ of maps inverted by
sheafification is a sameness, that $J\mapsto W_J$ is an order-embedding
$\Top(S)\hookrightarrow\Same(\Presh(S))$, that $W_J$ is the class of $J$-local
isomorphisms, and that the image consists exactly of the \emph{saturated}
samenesses that are the full class of maps inverted by a left-exact
reflector. Saturation is essential: $\{\mathrm{id}\}$ and $\Iso$ are distinct
samenesses with the same (identity) localization, and only the second is a
$W_J$.

\medskip\noindent\textbf{The finite-sieve dictionary.} To compute with sheaf
samenesses on the simplex and cube categories we use an explicit dictionary
(Theorem~\ref{thm:ideal}): on a category with finitely many sieves on each
object, Grothendieck topologies correspond, by reverse inclusion, to
\emph{idempotent two-sided ideals}, the ideal of a topology being the family
of its least covering sieves. This is the finite-sieve form of the
correspondence of Kelly and Lawvere \cite{KL89} between essential
localizations of a presheaf topos and idempotent ideals, and it is closely
related to the rigidity results of Johnstone \cite[C2.2.21]{Elephant},
Lindenhovius \cite{Lin14} and Marqu\`es \cite{Marques25}; we include the short
direct proof because it is elementary, needs no Cauchy completeness, and gives
the explicit least-cover description used later. Its consequences --- a finite
poset has $2^{|P|}$ topologies, a group two, a finite meet-semilattice $L$ as
a monoid has $\Top(L)\cong\Down(L^{\mathrm{op}})$ while $\Same(L)\cong
L^{\mathrm{op}}$, and the cube and simplex categories have exactly the
dimension chain of topologies, so that their subtoposes are the coskeletal
ones --- are re-derivations of, or small additions to, known facts: that the
essential subtoposes (levels) of simplicial and cubical sets are the
dimensions is due to Kennett, Riehl, Roy and Zaks \cite{KRRZ}, and that every
topology on $\Delta$ is rigid is in Marqu\`es \cite{Marques25}. What we add is
the position of these topologies as the sheaf chain $W_{D_0}\supseteq
W_{D_1}\supseteq\cdots$ in $\Same(\sset)$, which is what the main theorem
joins with $\Wtest$.

\medskip\noindent\textbf{Verification.} Every finite assertion (numbers of
topologies, the bijection with idempotent ideals, the poset, monoid,
semilattice and cube computations) was verified by enumeration from the
definition of a Grothendieck topology (\S\ref{sec:comp}).

\medskip\noindent\textbf{Related work.} Subtoposes of presheaf toposes and
their correspondence with Grothendieck topologies, local isomorphisms and
left-exact reflections are classical \cite{MLM,Elephant}; subtoposes of
$\Sh(X)$ for a locale $X$ correspond to nuclei \cite{Stone}. Kelly and
Lawvere \cite{KL89} showed that the essential localizations of a topos form a
complete lattice and described those of a presheaf topos $\Presh(S)$ through
idempotent two-sided ideals of $S$; Kennett, Riehl, Roy and Zaks \cite{KRRZ}
computed the resulting levels of simplicial, cubical and reflexive globular
sets and their Aufhebung; Marqu\`es \cite{Marques25} characterised, modulo
slices, the categories on which every Grothendieck topology is rigid, which
include finite Cauchy-complete categories, Artinian posets \cite{Lin14} and
$\Delta$. Toposes of monoid actions are studied by Hemelaer and Rogers
\cite{HR,Rogers}. Coskeleta, Kan's $\mathrm{Ex}^\infty$ and Postnikov sections
are in \cite{GJ}; that the coskeleta of a Kan complex model its Postnikov
sections is folklore, and we include the short proof
(Lemma~\ref{lem:cosk}). Composites and comparisons of localizations are a
classical theme of the theory of Bousfield localization and homotopy
idempotent functors \cite{Bousfield,Hirschhorn}, where localizations are
typically compared within model or homotopical structures and with
additional localization or saturation structure; the join theorem (Theorem~\ref{thm:generaljoin}) and
its corollaries assume none of this --- no model structure, no factorization,
no saturation --- and hold in the complete lattice of bare two-out-of-three
classes of an arbitrary category, which is what allows the sheaf-theoretic
class $W_{D_{n+1}}$ to be joined with $\Wtest$ without first realising both
as weak equivalences of model structures. Weak-equivalence
classes, localizers and model structures on presheaf categories are studied
by Cisinski and others; what is new here is the treatment of the bare
two-out-of-three classes as a lattice in which sheaf and homotopy
localizations can be joined, and the resulting identification of the
Postnikov truncations as such joins.

\section{Preliminaries}\label{sec:prelim}

\subsection{Samenesses}
For a small category $C$, a \emph{sameness} is a class $W\subseteq\Mor(C)$
with (i) all identities in $W$, (ii) $W$ closed under composition, (iii)
two-out-of-three: for composable $f,g$, if two of $f,g,g\circ f$ lie in $W$ then
so does the third. $\Same(C)$ is the set of samenesses ordered by inclusion.

\begin{prop}\label{prop:same}
$\Same(C)$ is a complete lattice (meets are intersections; $\bot=\{\mathrm{id}\}$,
$\top=\Mor(C)$). For a functor $F\colon C\to D$ and $W\in\Same(D)$,
$F^*(W)=\{m:F(m)\in W\}$ is a sameness in $C$, and $F^*$ preserves meets. In
particular $V_F:=F^*(\Iso D)$ is the largest sameness inverted by $F$; $F$
factors through the localization $C\to C[V_F^{-1}]$, and the induced functor
$C[V_F^{-1}]\to D$ is an equivalence iff $F$ exhibits $D$ as the localization
of $C$ at $V_F$. For a group $G$ (as a one-object category), $\Same(G)$ is the
subgroup lattice $\Sub(G)$.
\end{prop}
\begin{proof}
The three conditions are preserved by intersections and hold for $\Mor(C)$;
$F^*$ preserves them because $F$ preserves identities and composites, and it
preserves intersections. The universal property of the localization
\cite[I.1]{GZ} gives the factorization, and the last assertion is then the
definition of ``exhibits as the localization''. For a group, two-out-of-three
applied to $(w^{-1},w)$, whose composite is the identity, forces $w^{-1}\in W$,
so a sameness is a subgroup; conversely every subgroup satisfies the three
conditions.
\end{proof}

\begin{defi}\label{def:sat}
For $W\in\Same(C)$, an object $Z$ is \emph{$W$-local} if $\mathrm{Hom}(f,Z)$
is bijective for every $f\in W$; write $\mathrm{Loc}(W)$ for the class of
$W$-local objects. The \emph{saturation} of $W$ is
$\mathrm{Sat}(W):=\{f:\mathrm{Hom}(f,Z)\text{ bijective for all }Z\in
\mathrm{Loc}(W)\}$, a sameness containing $W$; $W$ is \emph{saturated} if
$W=\mathrm{Sat}(W)$.
\end{defi}

\begin{lem}\label{lem:loc}
For $U,V\in\Same(C)$, $\mathrm{Loc}(U\vee V)=\mathrm{Loc}(U)\cap\mathrm{Loc}(V)$;
hence $\mathrm{Sat}(U\vee V)=\mathrm{Sat}(\mathrm{Sat}(U)\vee\mathrm{Sat}(V))$.
If $L\colon C\to C$ is a reflector onto a full replete subcategory
$\mathcal E$ (with unit $\eta$), then $W_L:=L^*(\Iso)$ is saturated,
$\mathrm{Loc}(W_L)=\mathcal E$, and $C\to\mathcal E$ is the localization of
$C$ at $W_L$.
\end{lem}
\begin{proof}
$Z$ is local for $U\vee V$ iff for both $U$ and $V$: one direction is
trivial, and conversely $\{f:\mathrm{Hom}(f,Z)\text{ bijective}\}$ is a
sameness containing $U$ and $V$, hence the join. The second assertion follows
since $\mathrm{Loc}(\mathrm{Sat}(U))=\mathrm{Loc}(U)$. For a reflector, $f$ is
inverted by $L$ iff $\mathrm{Hom}(Lf,Z)$ is bijective for all $Z\in\mathcal
E$ (Yoneda in $\mathcal E$) iff $\mathrm{Hom}(f,Z)$ is bijective for all
$Z\in\mathcal E$; so $\mathcal E\subseteq\mathrm{Loc}(W_L)$ and
$W_L=\{f:\mathrm{Hom}(f,Z)\text{ bij.\ }\forall Z\in\mathcal E\}$. If $Z$ is
$W_L$-local then $\eta_Z\in W_L$ gives $\mathrm{Hom}(\eta_Z,Z)$ bijective, so
$\eta_Z$ has a retraction $r$ with $r\eta_Z=\mathrm{id}$; then $\eta_Z r$ and
$\mathrm{id}_{LZ}$ have the same composite with $\eta_Z$, and
$\mathrm{Hom}(\eta_Z,LZ)$ is bijective, so $\eta_Z$ is an isomorphism and
$Z\in\mathcal E$. Hence $\mathrm{Loc}(W_L)=\mathcal E$ and $W_L$ is saturated.
The localization statement is \cite[I.1.3]{GZ}.
\end{proof}

\subsection{Sieves, topologies, sheafification}
A \emph{sieve} on an object $c$ of $S$ is a set of morphisms with codomain $c$
closed under precomposition; for $h\colon d\to c$ and a sieve $R$ on $c$,
$h^*R=\{g: h\circ g\in R\}$. A \emph{Grothendieck topology} $J$ assigns to each
$c$ a set $J(c)$ of sieves such that the maximal sieve is in $J(c)$, $J$ is
stable ($R\in J(c)\Rightarrow h^*R\in J(d)$) and transitive (if $R\in J(c)$ and
$Q$ is a sieve on $c$ with $h^*Q\in J(d)$ for all $h\colon d\to c$ in $R$, then
$Q\in J(c)$) \cite{MLM}. Each $J(c)$ is then upward closed and closed under
finite intersection. $\Top(S)$ is the set of topologies ordered by $J\le J'$
iff $J(c)\subseteq J'(c)$ for all $c$; it is a complete lattice. We write
$a_J\colon\Presh(S)\to\Sh_J(S)$ for sheafification, a left-exact left adjoint to
the inclusion; $J\le J'$ implies $a_{J'}=a_{J'}\circ a_J$ up to isomorphism.
A morphism of presheaves is a \emph{$J$-local isomorphism} if it is $J$-locally
injective and $J$-locally surjective; $a_J f$ is an isomorphism iff $f$ is a
$J$-local isomorphism \cite[III.7]{MLM}, \cite[C2.2]{Elephant}.

\section{Topologies embed in the sameness lattice}\label{sec:embed}

\begin{thm}\label{thm:embed}
Let $S$ be a small category. For $J\in\Top(S)$ put
$W_J:=\{f\in\Mor\Presh(S): a_J f\text{ is an isomorphism}\}$. Then:
\begin{enumerate}
\item $W_J\in\Same(\Presh(S))$, and $W_J$ is the class of $J$-local
isomorphisms; it is stable under pullback and closed under colimits in the
arrow category;
\item $J\mapsto W_J$ is an order-embedding $\Top(S)\hookrightarrow
\Same(\Presh(S))$;
\item $W_J=V_{a_J}$ is the largest sameness inverted by $a_J$, and the induced
$\Presh(S)[W_J^{-1}]\to\Sh_J(S)$ is an equivalence;
\item the image of the embedding is the set of saturated samenesses $W$
that are the full class of maps inverted by a left-exact reflector onto a
full replete subcategory of $\Presh(S)$; equivalently, $\mathrm{Loc}(W)$ is a
subtopos and $W=\mathrm{Sat}(W)$.
\end{enumerate}
\end{thm}
\begin{proof}
(1) $W_J=a_J^*(\Iso)$ is a sameness by Proposition~\ref{prop:same}; the local
isomorphism description is the classical characterisation of the maps
inverted by sheafification; pullback stability because $a_J$ is left exact,
colimit closure because it is a left adjoint. (2) Monotone: $J\le J'$ gives
$a_{J'}=a_{J'}a_J$, so $a_Jf$ iso implies $a_{J'}f$ iso. Injective and
order-reflecting: the $W_J$-local objects are exactly the $J$-sheaves (for a
reflective localization the local objects form the reflective subcategory),
and a subtopos of $\Presh(S)$ determines its topology \cite[V.4]{MLM}, \cite[A4.3]{Elephant}, so $W_J\subseteq
W_{J'}$ forces $\Sh_{J'}\subseteq\Sh_J$, i.e.\ $J\le J'$. (3) By
Lemma~\ref{lem:loc}, since $a_J$ is a reflector. (4) If $W=L^*(\Iso)$ for a
left-exact reflector $L$ onto $\mathcal E\subseteq\Presh(S)$, then
$\mathcal E$ is a subtopos, hence $\mathcal E=\Sh_J(S)$ for a unique $J$
\cite[A4.3]{Elephant}, and $W=\{f:\mathrm{Hom}(f,Z)\text{ bij.\ }\forall
Z\in\mathcal E\}=W_J$ by Lemma~\ref{lem:loc}; conversely $W_J=a_J^*(\Iso)$ is
of this form and saturated. The condition ``saturated'' cannot be dropped:
$\{\mathrm{id}\}$ and $\Iso$ have the same localization but only $\Iso=W_{J_{\min}}$
(trivial topology) is in the image.
\end{proof}

\begin{rem}
$\Presh(S)$ is large; as in \cite{SameA} we fix universes $U\in V$, take $S$
$U$-small and form $\Same$ and the localizations in $V$. Nothing below
depends on this bookkeeping. Figure~\ref{fig:embed} summarises
Theorem~\ref{thm:embed}.
\end{rem}

\begin{figure}[ht]
\centering
\begin{tikzcd}[column sep=large,row sep=large]
\Top(S)\arrow[r,hook,"J\mapsto W_J"]\arrow[d,"J\mapsto\Sh_J(S)"'] &
\Same(\Presh(S))\arrow[d,"{W\mapsto\Presh(S)[W^{-1}]}"] \\
\{\text{subtoposes}\}^{\mathrm{op}}\arrow[r,hook] & \{\text{localizations}\}
\end{tikzcd}
\qquad
\begin{tikzcd}[column sep=large,row sep=large]
\Presh(S)\arrow[r,"a_J"]\arrow[d,"\text{loc.}"'] & \Sh_J(S) \\
\Presh(S)[W_J^{-1}]\arrow[ur,"\simeq"'] &
\end{tikzcd}
\caption{Top: the order-embedding of Theorem~\ref{thm:embed} over the
classical correspondence between topologies and subtoposes; the right-hand
column sends a sameness to its localization, and the image of the top row is
the set of saturated samenesses inverted by a left-exact reflector. Bottom: $W_J$ is the
largest sameness inverted by $a_J$, and the induced comparison is an
equivalence.}
\label{fig:embed}
\par\smallskip\noindent\textit{Alt text:} Above, a commutative square: the
lattice of topologies embeds into the sameness lattice along the top, and
both map down to subtoposes and localizations respectively. Below, a
triangle: sheafification from presheaves to sheaves factors through the
localization at the sameness of the topology, and the comparison functor is
an equivalence.
\end{figure}

\section{Joins of samenesses and localizations}\label{sec:general}

The main theorem rests on the following formal statement.

\begin{thm}[join theorem]\label{thm:generaljoin}
Let $U,V,W\in\Same(C)$ with $W\subseteq U\vee V$, and let $P\colon C\to C$ be
a functor with a natural transformation $\pi\colon\mathrm{id}_C\Rightarrow P$
such that
\begin{enumerate}
\item[(a)] $\pi_X\in U\vee V$ for every object $X$;
\item[(b)] $P(U)\subseteq W$ and $P(V)\subseteq W$.
\end{enumerate}
Then $U\vee V=P^*(W)=\{f:Pf\in W\}$.
\end{thm}
\begin{proof}
$P^*(W)$ is a sameness (Proposition~\ref{prop:same}) containing $U$ and $V$
by (b), hence containing $U\vee V$. Conversely let $f\colon X\to Y$ with
$Pf\in W\subseteq U\vee V$. Naturality gives $\pi_Y\circ f=Pf\circ\pi_X$; the
right-hand side lies in $U\vee V$ by (a), and $\pi_Y\in U\vee V$, so
two-out-of-three gives $f\in U\vee V$.
\end{proof}

\begin{cor}[two reflectors]\label{cor:tworeflectors}
Let $L_1,L_2$ be reflectors of $C$ onto full replete subcategories, with
$W_{L_i}=L_i^*(\Iso)$, and write $W_{L_2L_1}:=(L_2L_1)^*(\Iso)$ for the class
of maps inverted by the composite functor (which need not itself be a
reflector). Then
\[
W_{L_1}\vee W_{L_2}=W_{L_2L_1}\quad\Longleftrightarrow\quad
L_2L_1\text{ inverts }W_{L_2}.
\]
In particular the equality holds whenever $L_1$ preserves $L_2$-equivalences,
or whenever $L_2L_1\cong L_2$ (e.g.\ $\mathrm{Loc}(W_{L_2})\subseteq
\mathrm{Loc}(W_{L_1})$).
\end{cor}
\begin{proof}
Apply Theorem~\ref{thm:generaljoin} with $U=W_{L_1}$, $V=W_{L_2}$, $W=\Iso$,
$P=L_2L_1$ and $\pi_X=\eta^2_{L_1X}\circ\eta^1_X$, a composite of a
$W_{L_1}$-map and a $W_{L_2}$-map; (b) for $U$ holds because $L_2L_1$ inverts
$W_{L_1}$, and (b) for $V$ is the stated condition. Conversely, if the
equality holds then $W_{L_2}\subseteq W_{L_2L_1}$.
\end{proof}

\begin{prop}[joins of topologies]\label{prop:topjoin}
For Grothendieck topologies $J,J'$ on $S$, with $J\vee J'$ their join in
$\Top(S)$,
\[
W_{J\vee J'}=\mathrm{Sat}(W_J\vee W_{J'}),
\]
so $W_{J\vee J'}=W_J\vee W_{J'}$ iff the join $W_J\vee W_{J'}$ is saturated.
\end{prop}
\begin{proof}
A presheaf is a $(J\vee J')$-sheaf iff it is a $J$-sheaf and a $J'$-sheaf,
because the sieves for which a given presheaf (and all its pullbacks)
satisfies the sheaf condition form a topology \cite[C2.1]{Elephant}. By
Lemma~\ref{lem:loc}, $\mathrm{Loc}(W_J\vee W_{J'})=\Sh_J\cap\Sh_{J'}
=\Sh_{J\vee J'}=\mathrm{Loc}(W_{J\vee J'})$, and $W_{J\vee J'}$ is saturated,
so $W_{J\vee J'}=\mathrm{Sat}(W_J\vee W_{J'})$.
\end{proof}

\begin{ex}\label{ex:twochain}
Let $S$ be the poset $0<1$, so that $\Presh(S)$ is the category of maps of
sets $X(1)\to X(0)$. The two intermediate topologies (\S\ref{sec:ideal}) have
sheaf samenesses $W_0=\{f:f_0\text{ bijective}\}$ and $W_1=\{f:f_1\text{
bijective}\}$, and their join topology is the degenerate one, with
$W=\Mor$. Here $W_0\vee W_1=\Mor$: any $f=(f_1,f_0)\colon X\to Y$ factors as
$(X(1)\to X(0))\xrightarrow{(\mathrm{id},f_0)}(X(1)\to Y(0))
\xrightarrow{(f_1,\mathrm{id})}(Y(1)\to Y(0))$, a $W_1$-map followed by a
$W_0$-map. Whether $W_J\vee W_{J'}=W_{J\vee J'}$ holds for all pairs of
topologies on all sites we leave open; Proposition~\ref{prop:topjoin} reduces
it to the saturation of the join.
\end{ex}

\section{The finite-sieve dictionary}\label{sec:ideal}

A \emph{two-sided ideal} of $S$ is a class $D$ of morphisms closed under pre-
and post-composition with arbitrary morphisms of $S$; write $D(c)$ for its
members with codomain $c$ (a sieve on $c$). $D$ is \emph{idempotent} if
$D\circ D=D$, i.e.\ every member of $D$ is a composite of two members of $D$.
Let $\Idem(S)$ be the set of idempotent two-sided ideals.

Kelly and Lawvere \cite{KL89} describe the essential localizations of
$\Presh(S)$ --- those whose reflector has a left adjoint --- in terms of
idempotent two-sided ideals of $S$; when every object has finitely many
sieves, every topology has least covering sieves and the description becomes
the following elementary dictionary, which we prove directly since it is
short, requires no Cauchy completeness, and is what we compute with. (For
categories with finite slices \cite[C2.2.21]{Elephant}, \cite[Prop.~4.10]{KL89},
and for those in which every arrow factors as a split epimorphism followed by
a split monomorphism and every object has finitely many subobjects (Menni;
see criterion (C2) in \cite{Marques25}), every topology is moreover
\emph{rigid}: its sheaf topos is a presheaf topos on the full subcategory
of $J$-irreducible objects, i.e.\ the objects $c$ with $\mathrm{id}_c\in D_J$
below.)

\begin{thm}[ideal theorem]\label{thm:ideal}
Let $S$ have finitely many sieves on each object (e.g.\ $S$ finite). Then
\[
\Top(S)\ \cong\ \Idem(S)^{\mathrm{op}},
\]
\[
J\longmapsto D_J:=\Bigl(\bigcap J(c)\Bigr)_{c},\qquad
D\longmapsto J_D(c):=\{R\text{ a sieve on }c: R\supseteq D(c)\},
\]
mutually inverse order-reversing bijections. The trivial topology corresponds
to $D=\Mor(S)$ and the degenerate one (all sieves) to $D=\varnothing$.
\end{thm}
\begin{proof}
Let $J$ be a topology. Each $J(c)$ is upward closed and closed under finite
intersection, so with finitely many sieves $D_J(c):=\bigcap J(c)\in J(c)$ and
$J(c)=\{R\supseteq D_J(c)\}$. Stability of $D_J(c)$ along $h\colon d\to c$ gives
$h^*D_J(c)\in J(d)$, hence $h^*D_J(c)\supseteq D_J(d)$, i.e.\ $h\circ g\in
D_J(c)$ for every $g\in D_J(d)$: $D_J$ is closed under post-composition, and
being a family of sieves it is closed under pre-composition, so it is a
two-sided ideal. Transitivity applied to $R=D_J(c)$ and $Q=(D_J\circ D_J)(c)$:
for $h\in D_J(c)$, $h^*Q\supseteq D_J(d)$ since $h\circ g\in D_J\circ D_J$ for
$g\in D_J(d)$; hence $Q\in J(c)$, so $Q\supseteq D_J(c)$, i.e.\ $D_J\circ
D_J\supseteq D_J$; the other inclusion is trivial, so $D_J$ is idempotent.

Conversely let $D$ be an idempotent two-sided ideal. $J_D(c)$ contains the
maximal sieve. Stability: if $R\supseteq D(c)$ and $g\in D(d)$ then $h\circ
g\in D(c)\subseteq R$, so $h^*R\supseteq D(d)$. Transitivity: let $R\supseteq
D(c)$ and $Q$ a sieve with $h^*Q\supseteq D(d)$ for all $h\colon d\to c$ in $R$;
in particular for $h\in D(c)$, so $h\circ D(d)\subseteq Q$ for all $h\in D(c)$,
i.e.\ $(D\circ D)(c)\subseteq Q$, whence $D(c)\subseteq Q$ by idempotence and
$Q\in J_D(c)$. The two assignments are inverse ($D_{J_D}=D$ since $D(c)$ is the
least sieve containing $D(c)$; $J_{D_J}=J$ as shown) and reverse inclusion.
\end{proof}

The theorem was checked by enumerating topologies from the definition and
idempotent two-sided ideals independently, on sixty-one monoids (all of
order $\le 4$, cyclic groups, thirteen lattices as meet-monoids) and on ten
categories with several objects, parallel arrows and non-trivial
endomorphisms; the bijection and its order-reversal held in every case
(\S\ref{sec:comp}); see Figure~\ref{fig:ideal}.

\begin{figure}[ht]
\centering
\begin{tikzcd}[column sep=huge,row sep=large]
J\arrow[r,mapsto,"{D_J=(\bigcap J(c))_c}"] & D_J \\
J_D\arrow[u,equal,"{J_{D_J}=J}"] & D\arrow[l,mapsto,"{J_D(c)=\{R\supseteq D(c)\}}"]\arrow[u,equal,"{D_{J_D}=D}"']
\end{tikzcd}
\caption{The ideal theorem (Theorem~\ref{thm:ideal}): a topology is recorded
by its least covering sieves, which form an idempotent two-sided ideal, and
an idempotent two-sided ideal generates a topology by upward closure; the two
passages are inverse and reverse the order.}
\label{fig:ideal}
\par\smallskip\noindent\textit{Alt text:} A square of assignments. A
topology $J$ maps right to the ideal of its least covering sieves; an ideal
$D$ maps left to the topology of sieves containing $D$; the vertical
equalities record that the two round trips are identities.
\end{figure}

\begin{cor}[posets]\label{cor:poset}
For a finite poset $P$ (as a category), $\Top(P)\cong 2^{|P|}$, the Boolean
lattice on the underlying set: the topologies depend only on the number of
points, not on the order. (This also follows from rigidity: the topologies on
an Artinian poset correspond to its subsets \cite{Lin14}, \cite{Marques25}.)
\end{cor}
\begin{proof}
For $Y\subseteq P$ let $D_Y$ be the set of arrows $a\to b$ with $a\le y\le b$
for some $y\in Y$; it is a two-sided ideal, idempotent since $a\to b$ factors
as $(y\to b)\circ(a\to y)$ with both factors in $D_Y$, and $Y$ is recovered
from $D_Y$ as the set of objects whose identity it contains. Conversely let
$D$ be an idempotent two-sided ideal and $Y$ the set of objects with
$\mathrm{id}_y\in D$; then $D\supseteq D_Y$. For $a\to b$ in $D$, idempotence
gives $y_1$ with $a\le y_1\le b$ and $a\to y_1$, $y_1\to b$ in $D$; if $y_1=a$
or $y_1=b$ an identity of $D$ lies between $a$ and $b$; otherwise repeat with
$a\to y_1$. Since $P$ is finite the strictly shrinking intervals terminate,
producing $y\in Y$ with $a\le y\le b$, so $a\to b\in D_Y$. Hence $D=D_Y$, and
$Y\mapsto D_Y$ is an order-isomorphism $2^P\cong\Idem(P)$. Verified from the
definition of a topology on eight posets with $\le4$ points (chains,
antichains, $\mathrm{V}$, $\Lambda$, the square), in each case also matching
the number of nuclei on the frame $\Down(P)$.
\end{proof}

\begin{cor}[groups, groupoids, monoids]\label{cor:monoid}
For a monoid $M$ with finitely many right ideals, $\Top(M)$ is the poset of
idempotent two-sided ideals of $M$ under reverse inclusion. For a group (or a
connected groupoid) the only two-sided ideals are $\varnothing$ and everything,
so $\Top(G)$ is the two-element chain, while $\Same(G)=\Sub(G)$; for a
groupoid with a set $\pi_0$ of components, $\Top\cong 2^{\pi_0}$.
\end{cor}

\begin{cor}[semilattices]\label{cor:semilattice}
Let $L$ be a finite meet-semilattice with top element, viewed as a monoid under
$\wedge$. Then every ideal (down-set) of $L$ is a two-sided idempotent ideal, so
\[
\Top(L)\ \cong\ \Down(L)^{\mathrm{op}}\ \cong\ \Down(L^{\mathrm{op}}),
\]
the lattice of up-sets of $L$; in particular the lattice of subtoposes of the
topos of $L$-sets determines $L$ up to isomorphism (Birkhoff duality). The
arrow-level sameness lattice is $\Same(L)\cong L^{\mathrm{op}}$ via $a\mapsto
{\uparrow}a$; hence $\Top(L)\cong\Down(\Same(L))$.
\end{cor}
\begin{proof}
A down-set $D$ satisfies $D\wedge L\subseteq D$ and $d=d\wedge d$, so it is a
two-sided idempotent ideal; conversely a right ideal is a down-set. For
$\Same(L)$: a sameness $W$ is a meet-closed submonoid; with $a:=\bigwedge W\in W$
we have $W\subseteq{\uparrow}a$, and for $y\ge a$ the pair $(y,a)$ has
composite $y\wedge a=a\in W$, so two-out-of-three gives $y\in W$; each
${\uparrow}a$ is a sameness, and ${\uparrow}a\subseteq{\uparrow}b$ iff $b\le a$.
Verified on thirteen lattices ($n$-chains, $2^2$, $2^3$, $M_3$, $N_5$, $2\times3$,
$1\oplus2^2$, $2^2\oplus1$, $2^2\oplus2$): e.g.\ $|\Top(2^3)|=20$, $|\Top(M_3)|=10$,
$|\Top(N_5)|=8$.
\end{proof}

\begin{ex}
Among the $35$ monoids of order $4$, $\Top(M)$ is a chain for $34$; the unique
exception is the meet-semilattice $2^2$ (identity, zero, two orthogonal
idempotents), with $\Top(2^2)\cong\Down(2^2)^{\mathrm{op}}$, six elements
(Figure~\ref{fig:semilattice}). The pairs $(|\Top M|,|\Same M|)$ take nine
different values on order $4$; neither invariant determines the other.
\end{ex}

\begin{figure}[ht]
\centering
\begin{tikzcd}[row sep=small, column sep=small]
& \top & \\
a\arrow[ur,dash] & & b\arrow[ul,dash] \\
& 0\arrow[ul,dash]\arrow[ur,dash] &
\end{tikzcd}
\qquad\qquad
\begin{tikzcd}[row sep=small, column sep=small]
& L & \\
& \{0,a,b\}\arrow[u,dash] & \\
\{0,a\}\arrow[ur,dash] & & \{0,b\}\arrow[ul,dash] \\
& \{0\}\arrow[ul,dash]\arrow[ur,dash] & \\
& \varnothing\arrow[u,dash] &
\end{tikzcd}
\caption{Left: the meet-semilattice $L=2^2$ as a monoid. Right: its six
topologies, $\Top(L)\cong\Down(L)^{\mathrm{op}}$ (drawn as $\Down(L)$; the
topology order is the reverse). The arrow-level lattice is $\Same(L)\cong
L^{\mathrm{op}}$, four elements.}
\label{fig:semilattice}
\par\smallskip\noindent\textit{Alt text:} Two Hasse diagrams. Left, a diamond
with bottom $0$, two incomparable elements $a$ and $b$, and top. Right, a
six-element lattice: empty set at the bottom, then the singleton of $0$, then
two incomparable sets $\{0,a\}$ and $\{0,b\}$, then $\{0,a,b\}$, then the whole
semilattice at the top.
\end{figure}

\section{Cube and simplex categories}\label{sec:cube}

Let $\cube$ be the cube category generated by faces and degeneracies (Kan's
cube category, as in \cite{KRRZ}; no connections, reversals or symmetries)
and $\Delta$ the simplex category. For $k\ge0$ let $D_k$ be the class of
morphisms that factor through an object of dimension $\le k$. The following
theorem is a re-derivation, through the dictionary, of known facts: that the
levels (essential subtoposes) of simplicial and cubical sets are the
dimensions is \cite[\S1]{KRRZ}, and that every topology on $\Delta$ is rigid
is \cite[Thm.~2.6]{Marques25}, so that all subtoposes of $\sset$ are levels;
the same holds for $\cube$, which satisfies the split-epi/split-mono
factorization criterion of Menni recorded as (C2) in \cite{Marques25}. What we need from it is
the last sentence: the explicit description of the sheaf samenesses.

\begin{thm}\label{thm:cube}
The two-sided ideals of $\cube$ (resp.\ $\Delta$) are exactly $\varnothing$,
the $D_k$ ($k\ge0$) and $\Mor$; all are idempotent. Hence the topologies form
the chain
\[
\text{degenerate}\ \subset\ \cdots\ \subset\ J_{D_1}\ \subset\ J_{D_0}\ \subset\ \text{trivial}
\]
(reverse inclusion of ideals), and the $J_{D_k}$-sheaves are the
$k$-coskeletal presheaves: the subtoposes of cubical sets and of simplicial
sets are precisely the coskeletal ones. The sheaf samenesses form the
descending chain $W_{D_0}\supseteq W_{D_1}\supseteq\cdots$ of
$k$-coskeletal equivalences,
\[
W_{D_k}=\{f:\cosk_kf\text{ iso}\}=\{f: f_m\text{ bijective for all }m\le k\},
\]
with $\bigcap_kW_{D_k}=\Iso$.
\end{thm}
\begin{proof}
Every morphism $f\colon[m]\to[n]$ factors as a degeneracy $\sigma\colon[m]\twoheadrightarrow[k]$
followed by a face $\delta\colon[k]\hookrightarrow[n]$, $k$ the rank of $f$;
degeneracies have sections and faces have retractions among the generators.
So $\mathrm{id}_{[k]}=r\circ f\circ s$ lies in the two-sided ideal generated by
$f$, which therefore contains every morphism factoring through $[k]$, i.e.\
$D_k$, and $f\in D_k$. Hence principal ideals are the $D_k$, two-sided ideals
are unions of $D_k$'s --- $\varnothing$, some $D_k$, or everything --- and
each is idempotent since it contains an identity through which its members
factor. Theorem~\ref{thm:ideal} gives the topologies: its finiteness hypothesis
holds, a sieve on $[n]$ being a subpresheaf of the representable $y[n]$,
which is determined by the finitely many non-degenerate cells it contains
(every cell of $y[n]$ is uniquely a degeneracy of a non-degenerate one). The least
covering sieve of $J_{D_k}$ on $[n]$ is $D_k([n])$, generated by the maps from
dimension $\le k$, so a sheaf $X$ satisfies $X([n])\cong\mathrm{Hom}(\sk_ky[n],X)$ for all
$n$, i.e.\ $X\cong\cosk_kX$; conversely a $k$-coskeletal $X$ is local for
every sieve $R\supseteq\sk_ky[n]$, since $\mathrm{Hom}(R,\cosk_kX)=
\mathrm{Hom}(\sk_kR,X)$ and $\sk_kR=\sk_ky[n]$. Since $(\cosk_kX)_m=\mathrm{Hom}(\sk_ky[m],X)$
depends only on the $k$-truncation of $X$ and equals $X_m$ for $m\le k$,
$\cosk_kf$ is an isomorphism iff $f$ is bijective in dimensions $\le k$; the
intersection is therefore $\Iso$.
\end{proof}
Verified on the truncations $\cube_{\le1}$ ($7$ morphisms), $\cube_{\le2}$
($37$), $\Delta_{\le1}$ ($7$), $\Delta_{\le2}$ ($31$): exactly
$\varnothing,D_0,\dots,D_N$ in each case. Figure~\ref{fig:chain} shows the
chain together with its sheaves.

\begin{figure}[ht]
\centering
\begin{tikzcd}[row sep=small,column sep=small]
\text{ideal } D & \text{topology } J_D & \text{sheaves} & \text{sameness } W_D\\
\varnothing & \text{degenerate} & \text{terminal only} & \Mor \\
D_0\arrow[u,dash] & J_{D_0}\arrow[u,dash] & 0\text{-coskeletal} & W_{D_0}\arrow[u,dash]\\
D_1\arrow[u,dash] & J_{D_1}\arrow[u,dash] & 1\text{-coskeletal} & W_{D_1}\arrow[u,dash]\\
\vdots\arrow[u,dash] & \vdots\arrow[u,dash] & \vdots & \vdots\arrow[u,dash]\\
\Mor\arrow[u,dash] & \text{trivial}\arrow[u,dash] & \text{all presheaves} & \Iso\arrow[u,dash]
\end{tikzcd}
\caption{Theorem~\ref{thm:cube}: the four faces of the dimension chain for
$\cube$ and $\Delta$. Ideals grow downwards, topologies and samenesses grow
upwards; the sheaves for $J_{D_k}$ are the $k$-coskeletal presheaves.}
\label{fig:chain}
\par\smallskip\noindent\textit{Alt text:} A table of four columns aligned as
a vertical chain: the ideals from empty to everything, the corresponding
topologies from degenerate to trivial, their sheaves from the terminal object
through the coskeletal presheaves to all presheaves, and the samenesses from
all morphisms down to the isomorphisms.
\end{figure}

\section{Sheaf sameness versus homotopy sameness}\label{sec:join}

On $\sset=\Presh(\Delta)$ the lattice $\Same(\sset)$ contains the homotopy
sameness $\Wtest$ (weak homotopy equivalences) and the sheaf chain $W_{D_k}$.
Let $\mathrm{Ex}^\infty\colon\sset\to\sset$ be Kan's fibrant replacement, with
its natural weak equivalence $X\to\mathrm{Ex}^\infty X$ into a Kan complex
\cite[III.4]{GJ}. For $n\ge-1$ put
\[
P_n:=\cosk_{n+1}\circ\mathrm{Ex}^\infty\colon\sset\to\sset,
\]
\[
T_n:=P_n^*(\Wtest)=\{f: P_nf\text{ is a weak equivalence}\}.
\]
By Proposition~\ref{prop:same}, $T_n$ is a sameness. The following lemma is
standard: the coskeleta of a Kan complex model its Postnikov tower (Duskin
\cite{Duskin}; cf.\ the Postnikov towers of \cite[VI]{GJ}). We include the
elementary proof for completeness, since the main theorem rests on it.

\begin{lem}\label{lem:cosk}
Let $X$ be a Kan complex and $n\ge-1$, and put $K=\cosk_{n+1}X$. Then $K$ is
a Kan complex; the unit $\eta\colon X\to K$ induces isomorphisms
$\pi_i(X,x)\to\pi_i(K,x)$ for all $0\le i\le n$ and all vertices $x$; and
$\pi_i(K,x)=0$ for all $i\ge n+1$ and all $x$ (for $n=-1$: $K$ is empty or
contractible). Consequently, for $n\ge0$,
\[
T_n=\{f:\ \pi_0f\text{ bijective, }\pi_if\text{ iso at every base point, }1\le i\le n\},
\]
and $T_{-1}$ is the class of maps $X\to Y$ with $X=\varnothing$ iff
$Y=\varnothing$.
\end{lem}
\begin{proof}
An $m$-simplex of $K$ is a map $\sk_{n+1}\Delta^m\to X$, and $\eta$ is
bijective in dimensions $\le n+1$. \emph{Kan condition.} A horn
$\Lambda^k_m\to K$ is a map $\sk_{n+1}\Lambda^k_m\to X$. If $m\ge n+3$ then
$\sk_{n+1}\Lambda^k_m=\sk_{n+1}\Delta^m$, so the horn is already an
$m$-simplex of $K$. If $m\le n+2$ then $\Lambda^k_m$ has dimension $\le n+1$,
so the horn is a horn in $X$; a filler $\Delta^m\to X$ restricts to
$\sk_{n+1}\Delta^m\to X$, an $m$-simplex of $K$ extending it.
\emph{Homotopy groups.} For a Kan complex, $\pi_i(-,x)$ is computed from the
$i$-simplices with boundary at $x$ modulo the relation given by
$(i{+}1)$-simplices \cite[I.7]{GJ}; for $i\le n$ these dimensions are
$\le n+1$, where $\eta$ is bijective, so $\pi_i\eta$ is an isomorphism.
For $i\ge n+1$, a map $\partial\Delta^{i+1}\to K$ is a map
$\sk_{n+1}\partial\Delta^{i+1}=\sk_{n+1}\Delta^{i+1}\to X$, i.e.\ an
$(i{+}1)$-simplex of $K$ extending it; a Kan complex with this extension
property for all $i\ge n+1$ has $\pi_i=0$ for $i\ge n+1$ at all base points
\cite[I.7]{GJ}. \emph{The description of $T_n$.} A map of Kan complexes is a
weak equivalence iff it induces isomorphisms on $\pi_0$ and on all $\pi_i$
at all base points, and $\mathrm{Ex}^\infty$ preserves homotopy groups
\cite[III.4]{GJ}; apply this to $\cosk_{n+1}\mathrm{Ex}^\infty f$ using the
two previous assertions.
\end{proof}

\begin{thm}[join $=$ truncation]\label{thm:join}
In $\Same(\sset)$, for every $n\ge-1$,
\[
\Wtest\vee W_{D_{n+1}}\ =\ T_n .
\]
\end{thm}
\begin{proof}
We apply Theorem~\ref{thm:generaljoin} with $U=\Wtest$, $V=W_{D_{n+1}}$,
$W=\Wtest$ and $P=P_n$, so that $P^*(W)=T_n$ by definition. The natural
transformation $\pi_X\colon X\to\mathrm{Ex}^\infty X\to\cosk_{n+1}\mathrm{Ex}^\infty X$
is the Kan replacement, a weak equivalence, followed by the coskeleton unit,
which is bijective in dimensions $\le n+1$ (Theorem~\ref{thm:cube}); so
$\pi_X\in V\circ U\subseteq U\vee V$, which is (a). For (b): $P_n$ preserves
weak equivalences. Indeed $\mathrm{Ex}^\infty$ does; and if $g\colon X\to Y$ is
a weak equivalence between Kan complexes, then by Lemma~\ref{lem:cosk} the
units $X\to\cosk_{n+1}X$ and $Y\to\cosk_{n+1}Y$ induce isomorphisms on
$\pi_i$ for $i\le n$, so naturality shows that $\cosk_{n+1}g$ induces
isomorphisms on $\pi_i$ for $i\le n$ at all base points, while for $i>n$
both $\pi_i(\cosk_{n+1}X)$ and $\pi_i(\cosk_{n+1}Y)$ vanish; as
$\cosk_{n+1}X$ and $\cosk_{n+1}Y$ are Kan, $\cosk_{n+1}g$ is a weak
equivalence. Thus $P_n(U)\subseteq W$. If $f\in W_{D_{n+1}}$ then $f$ is
bijective in dimensions $\le n+1$, so $|\sk_{n+1}X|\to|\sk_{n+1}Y|$ is a
homeomorphism. Now $|\sk_{n+1}X|$ is the $(n{+}1)$-skeleton of the CW complex
$|X|$ (one cell per non-degenerate simplex), and the inclusion of the
$(n{+}1)$-skeleton of a CW complex induces isomorphisms on $\pi_i$ for
$0\le i\le n$ at every base point (the pair $(|X|,|\sk_{n+1}X|)$ is
$(n{+}1)$-connected \cite[Cor.~4.12]{Hatcher}). In the commutative square of
$|\sk_{n+1}f|$ and $|f|$ with these inclusions, three maps induce
isomorphisms on $\pi_i$, $i\le n$, hence so does $|f|$; and $f$ preserves
emptiness. Hence $f\in T_n$ by Lemma~\ref{lem:cosk}, i.e.\ $P_nf\in W$,
which is $P_n(V)\subseteq W$. Figure~\ref{fig:join} displays the naturality square used in the proof of
Theorem~\ref{thm:generaljoin} in this case.
\end{proof}

\begin{figure}[ht]
\centering
\begin{tikzcd}[column sep=large,row sep=large]
X\arrow[r,"f"]\arrow[d,"\Wtest"'] & Y\arrow[d,"\Wtest"] \\
X'\arrow[r,"f'"]\arrow[d,"\eta_{X'}"',"W_{D_{n+1}}"] & Y'\arrow[d,"\eta_{Y'}","W_{D_{n+1}}"'] \\
\cosk_{n+1}X'\arrow[r,"P_nf"',"\Wtest"] & \cosk_{n+1}Y'
\end{tikzcd}
\caption{Theorem~\ref{thm:join} as an instance of
Theorem~\ref{thm:generaljoin}: the vertical maps of the upper square are
weak equivalences (Kan replacement), those of the lower square are
coskeletal equivalences, and the bottom map is a weak equivalence when
$f\in T_n$; two-out-of-three moves membership in the join from the bottom row
to $f$.}
\label{fig:join}
\par\smallskip\noindent\textit{Alt text:} Two stacked commutative squares.
Top row: $f$ from $X$ to $Y$. Middle row: the Kan replacement $f'$ from $X'$
to $Y'$, reached by vertical weak equivalences. Bottom row: the coskeleton of
$f'$, labelled as a weak equivalence, reached by vertical coskeletal
equivalences.
\end{figure}

So the Postnikov truncations --- at the $\infty$-level the $n$-truncation
modalities --- are exactly the joins, in the sameness lattice of simplicial
sets, of the homotopy sameness with the sheaf samenesses of the coskeletal
topologies; at $n=-1$ the join with the $0$-coskeletal sameness is the
$(-1)$-truncation, which identifies all non-empty simplicial sets. This is a
comparison of two theories carried out entirely by a lattice operation, and
it needs no fibrancy hypothesis: the Kan replacement is itself available
inside $\Wtest$. Restricting along $\Kan\hookrightarrow\sset$
(Proposition~\ref{prop:same}) the same identity holds in $\Same(\Kan)$.

\begin{cor}[sheaf-theoretic Whitehead]\label{cor:whitehead}
In $\Same(\sset)$,
\[
\Wtest=\bigcap_{n\ge0}\bigl(\Wtest\vee W_{D_n}\bigr),\qquad
\Iso=\bigcap_{n\ge0}W_{D_n}.
\]
\end{cor}
\begin{proof}
By Theorem~\ref{thm:join} the first intersection is $\bigcap_{n\ge-1}T_n$,
the maps inducing isomorphisms on all homotopy groups at all base points,
which is $\Wtest$ by Whitehead's theorem (after Kan replacement). The second
is Theorem~\ref{thm:cube}.
\end{proof}

Thus the homotopy sameness is the meet of its joins with the sheaf chain,
while the sheaf chain itself meets down to the isomorphisms: the two families
are ``transversal'' in $\Same(\sset)$, and the truncations $T_n$ are the
interval they span (Figure~\ref{fig:lattice}).

\begin{figure}[ht]
\centering
\begin{tikzcd}[row sep=small,column sep=small]
& \Mor & \\
& T_{-1}=\Wtest\vee W_{D_0}\arrow[u,dash] & \\
& T_{0}=\Wtest\vee W_{D_1}\arrow[u,dash] & \\
& T_{n}=\Wtest\vee W_{D_{n+1}}\arrow[u,dash,dotted] & \\
\Wtest\arrow[ur,dash] & & W_{D_{n+1}}\arrow[ul,dash] \\
& \Iso=\bigcap_kW_{D_k}\arrow[ul,dash]\arrow[ur,dash] &
\end{tikzcd}
\caption{The position of the sheaf chain and the homotopy sameness in
$\Same(\sset)$ (Theorem~\ref{thm:join} and Corollary~\ref{cor:whitehead}):
the joins along the chain are the Postnikov truncations, decreasing to
$\Wtest$; the chain itself decreases to $\Iso$.}
\label{fig:lattice}
\par\smallskip\noindent\textit{Alt text:} A Hasse-style diagram. At the
bottom the isomorphisms; above them, to the left the weak equivalences and
to the right the coskeletal sameness of level $n+1$; their join, the
$n$-truncated equivalences, sits above both and is part of a chain of
truncations rising to the $(-1)$-truncation and then to all morphisms.
\end{figure} The analogous statement for cubical sets, where the
$n$-coskeletality of \cite{Cubes} concerns uniformly fibrant objects, is the
natural next step.

\section{Computational record}\label{sec:comp}
All finite claims are verified from the definitions by the Python scripts
supplied as supplementary material:
\begin{itemize}
\item \texttt{sheaf\_sameness.py}: topologies on posets from the definition,
nuclei on $\Down(P)$ (Corollary~\ref{cor:poset});
\item \texttt{sheaf\_sameness\_monoid.py}, \texttt{sheaf\_sameness\_monoid4.py}:
monoids of order $\le3$ and all $35$ monoids of order $4$;
\item \texttt{sheaf\_sameness\_semilattice.py}: thirteen lattices as
meet-monoids (Corollary~\ref{cor:semilattice});
\item \texttt{sheaf\_sameness\_ideals.py}: Theorem~\ref{thm:ideal} on $61$
monoids; \texttt{sheaf\_sameness\_general.py}: Theorem~\ref{thm:ideal} on ten
categories with several objects;
\item \texttt{sheaf\_sameness\_cube.py}: Theorem~\ref{thm:cube} on the
truncations $\cube_{\le2}$ and $\Delta_{\le2}$.
\end{itemize}

\begin{rem}[the sheaf level-shift logic]
In \cite{SameB} the modal logic of the frame $(\Same(C),\subseteq)$ was
studied. Restricting to the sheaf samenesses gives the modal logic of the lattice
of subtoposes, which by Theorem~\ref{thm:ideal} is, for a category with
finitely many sieves per object, the logic of its poset of idempotent
two-sided ideals: a combinatorial invariant of composition, computable without
forming a presheaf. For $\cube$ and $\Delta$ this frame is the chain of
Theorem~\ref{thm:cube}, so the logic contains $\mathrm{S4.3}$; for a finite
poset on $n$ points it is the logic of the Boolean frame $2^n$; for every
monoid of order $\le4$ except the semilattice $2^2$ it is the logic of a
finite chain.
\end{rem}

\section{Conclusion}
Sheaf theory's notion of sameness is a family of saturated elements of the
sameness lattice, intrinsically described (local isomorphisms), and within
these elements the join of two topologies is the saturation of the join of the
corresponding samenesses. On the cube and simplex categories the sheaf samenesses form the
dimension chain of coskeletal equivalences, and on simplicial sets the join
of this chain with the homotopy sameness is exactly the Postnikov truncation,
for every $n\ge-1$, with the homotopy sameness recovered as the meet of these
joins; the mechanism is a general join theorem for samenesses. Two questions are natural next steps:
the cubical version of Theorem~\ref{thm:join}, where the $n$-coskeletality of
\cite{Cubes} is precisely the sheaf sameness; and the $\infty$-categorical
statement, where reflective subuniverses (modalities) replace reflective
samenesses and Theorem~\ref{thm:join} should read ``truncation $=$ homotopy
$\vee$ coskeleton'' in the lattice of modalities.

\section*{Funding}
No funding was received for this work.

\section*{Declarations}
\noindent\textbf{Competing interests.}\quad The author declares that he has no
competing interests.

\noindent\textbf{Corresponding author.}\quad Ryota Shibusawa, Daiichi Institute
of Technology; e-mail: \texttt{r-shibusawa@daiichi-koudai.ac.jp}.

\noindent\textbf{Availability of data and materials.}\quad All verification
scripts named in \S\ref{sec:comp} are provided as supplementary material.

\begin{thebibliography}{9}
\bibitem{SameA} R.~Shibusawa, \emph{The sameness lattice of a category:
weak-equivalence structures, Galois functoriality, and comparison intervals},
preprint (2026), archived at DOI 10.5281/zenodo.22823772.
\bibitem{SameB} R.~Shibusawa, \emph{The level-shift logic of a category: modal
correspondence for sameness lattices, quotient models with contingent
identity, and completeness}, manuscript, under review (2026).
\bibitem{Cubes} R.~Shibusawa, \emph{A complexity dichotomy for order-convex
monotone selection, with applications to Dedekind cubical sets}, manuscript,
under review (2026).
\bibitem{MLM} S.~Mac Lane, I.~Moerdijk, \emph{Sheaves in Geometry and Logic},
Universitext, Springer, 1992.
\bibitem{Elephant} P.~T.~Johnstone, \emph{Sketches of an Elephant: A Topos
Theory Compendium}, Oxford Logic Guides \textbf{43--44}, Oxford University
Press, 2002.
\bibitem{Stone} P.~T.~Johnstone, \emph{Stone Spaces}, Cambridge Studies in
Advanced Mathematics \textbf{3}, Cambridge University Press, 1982.
\bibitem{GJ} P.~G.~Goerss, J.~F.~Jardine, \emph{Simplicial Homotopy Theory},
Progress in Mathematics \textbf{174}, Birkh\"auser, 1999.
\bibitem{Bousfield} A.~K.~Bousfield, \emph{The localization of spaces with
respect to homology}, Topology \textbf{14} (1975), 133--150.
\bibitem{Duskin} J.~Duskin, \emph{Simplicial methods and the interpretation
of ``triple'' cohomology}, Mem. Amer. Math. Soc. \textbf{163} (1975).
\bibitem{Hatcher} A.~Hatcher, \emph{Algebraic Topology}, Cambridge University
Press, 2002.
\bibitem{Hirschhorn} P.~S.~Hirschhorn, \emph{Model Categories and Their
Localizations}, Math. Surveys Monogr. \textbf{99}, Amer. Math. Soc., 2003.
\bibitem{GZ} P.~Gabriel, M.~Zisman, \emph{Calculus of Fractions and Homotopy
Theory}, Ergebnisse der Mathematik \textbf{35}, Springer, 1967.
\bibitem{KL89} G.~M.~Kelly, F.~W.~Lawvere, \emph{On the complete lattice of
essential localizations}, Bull. Soc. Math. Belg. S\'er. A \textbf{41} (1989),
289--319.
\bibitem{KRRZ} C.~Kennett, E.~Riehl, M.~Roy, M.~Zaks, \emph{Levels in the
toposes of simplicial sets and cubical sets}, J. Pure Appl. Algebra
\textbf{215} (2011), 949--961.
\bibitem{Lin14} B.~Lindenhovius, \emph{Grothendieck topologies on posets},
arXiv:1405.4408 (2014).
\bibitem{Marques25} J.~Marqu\`es, \emph{A criterion for categories on which
every Grothendieck topology is rigid}, Appl. Categ. Structures \textbf{33}
(2025), art.\ 37.
\bibitem{HR} J.~Hemelaer, M.~Rogers, \emph{Monoid properties as invariants of
toposes of monoid actions}, Appl. Categ. Structures \textbf{29} (2021),
379--413.
\bibitem{Rogers} M.~Rogers, \emph{Toposes of discrete monoid actions},
arXiv:1905.10277 (2019).
\end{thebibliography}
\end{document}
